% Generated draft: requires mathematical and source review.
\documentclass[11pt]{article}
\usepackage[margin=25mm]{geometry}
\usepackage{fontspec}
\setmainfont{Latin Modern Roman}
\usepackage{amsmath,amssymb,mathtools,unicode-math}
\setmathfont{Latin Modern Math}
\usepackage{enumitem,needspace,xurl,hyperref}
\setlength{\parindent}{0pt}
\setlength{\parskip}{5pt}
\setlength{\emergencystretch}{3em}
\newcommand{\N}{\mathbb{N}}
\newcommand{\Z}{\mathbb{Z}}
\newcommand{\Q}{\mathbb{Q}}
\newcommand{\R}{\mathbb{R}}
\newcommand{\C}{\mathbb{C}}
\newcommand{\F}{\mathbb{F}}
\newcommand{\A}{\mathbb{A}}
\newcommand{\PP}{\mathbb{P}}
\newcommand{\E}{\mathbb{E}}
\newcommand{\eps}{\varepsilon}
\newcommand{\dd}{\,\mathrm{d}}
\begin{document}
\section*{UnsolvedMath ID: 30003035}
\textbf{Source problem code:} catalogue entry
\subsection{Definitions and mathematical statement}
A conventional Euclidean formulation of the formal-duality property is as follows. Let \(n\geq 1\), let \(\Lambda\subset\R^n\) be a full-rank lattice, and let
\[
\mathcal P=\bigcup_{j=1}^{N}(x_j+\Lambda),
\]
where \(N\geq1\) and the representatives \(x_j\) are distinct modulo \(\Lambda\). Its density is
\[
\rho(\mathcal P)=\frac{N}{\operatorname{covol}(\Lambda)}.
\]
Let \(\mathcal S(\R^n)\) denote the space of smooth complex-valued functions whose derivatives decay faster than every inverse power of the Euclidean norm. For \(f\in\mathcal S(\R^n)\), use the Fourier-transform convention
\[
\widehat f(y)=\int_{\R^n}f(x)e^{-2\pi i\langle x,y\rangle}\,\dd x
\]
and define the averaged pair sum
\[
\Sigma_f(\mathcal P)=\frac1N\sum_{j=1}^{N}\sum_{k=1}^{N}\sum_{\lambda\in\Lambda}
 f(x_j-x_k+\lambda).
\]
This sum includes the self-pair terms \(j=k\), \(\lambda=0\); it is absolutely convergent and independent of the chosen lattice-period representation. For another periodic point set \(\mathcal Q\subset\R^n\), define \(\Sigma_g(\mathcal Q)\) in the same way. The formal-duality identity, with these conventions, is
\[
\Sigma_f(\mathcal P)=\rho(\mathcal P)\Sigma_{\widehat f}(\mathcal Q)
\qquad\text{for every }f\in\mathcal S(\R^n).
\]
Here \(\mathcal Q\) is a formal-dual partner; the definition does not require \(\mathcal Q=\mathcal P\).

\textbf{Source question.} Why do energy-minimizing periodic point sets appear to possess this formal-duality property?

\textbf{Unresolved qualifications.} The displayed definitions give a conventional precise meaning to the formal-duality identity, but the supplied record does not specify a dimension, an interaction potential or energy functional, density constraints, admissible competitors, the sense of minimization, or the numerical examples under discussion. It also does not specify how a formal-dual partner is selected or whether self-duality is intended in particular examples. Thus the source poses an explanatory question about an apparent phenomenon, rather than a fully quantified assertion that every periodic energy minimizer has a formal-dual partner. No such universal assertion is imposed here.
\subsection{Short English statement}
The problem asks why periodic configurations observed to minimize energy seem to admit a partner whose Fourier-transformed averaged pair sums reproduce their own pair sums after multiplication by density. The formal-duality identity can be stated precisely, but the supplied record leaves the energy model, constraints, and scope of the observation unspecified.
\subsection{Sources}
\begin{itemize}
\item Local UnsolvedMath registry: supplied complete source record for permanent ID \(30003035\), titled \emph{Formal Duality of Energy Minimizing Periodic Point Sets}; exact source problem code \texttt{catalogue entry}. The record assigns the year 2016 and names no individual proposer. Its statement status is supplied as \texttt{exact}, with the qualification that the wording was retained byte-for-byte after an internal readability check.
\item Supplied report attribution: \emph{Lattices and Applications in Number Theory} (2016), Oberwolfach Reports, pp.\ 87--154, DOI \texttt{10.4171/owr/2016/3}. \url{https://doi.org/10.4171/owr/2016/3}. The supplied background describes Oberwolfach Reports as documenting mathematical workshops and their open-problem discussions. The literature-source label for this same link is \emph{Oberwolfach Report 3/2016}.
\item Supplied literature reference: \emph{Formal duality in finite abelian groups}. \url{https://doi.org/10.1016/j.jcta.2018.11.005}.
\item Supplied dated literature curation, checked \texttt{2026-08-21}: the record labels the problem open and reports that formal duality has been formalized and partially classified in finite abelian groups, while no general theorem explains the numerically observed phenomenon. This is supplied dated curation derived from the Oberwolfach curation, not an independently verified literature assessment. The supplied background explicitly calls for independent verification of time-sensitive status and cited claims.
\end{itemize}
\end{document}
